\section{Proof of the main theorem} 
\label{section:3_morphism}

\subsubsection{} In this final section we will prove the main theorem in the following form:

\begin{theorem}
\label{thm:theorem_Ran}
The natural transformation $\eta_{G,\Ran}$ is an isomorphism.
\end{theorem}

\subsection{Proof of Theorem \ref{thm:theorem_Ran}}

\subsubsection{} Our proof proceeds by induction on $G$. We make a series of reductions.

\subsubsection{First reduction} We already know that $\eta_{G,\Ran}$ has Eisenstein cofiber. Therefore, it suffices to show that the cofiber of $\eta_{G, \Ran}$ is cuspidal. By Hecke linearity, it suffices to show that the cofiber of
    \[\eta_G \simeq \eta_{G,\Ran}(\WW_\vac)\]
is cuspidal.

\subsubsection{Second reduction} Fix a point $x \in X$ and choose $\lambda$ so that
    \[\Bun_N^{\lambda \cdot x} = \Bun_B \underset{\Bun_T}{\times} \{\omega^{\rhov}(\lambda \cdot x)\}\]
is a scheme. Write
    \[\WW_\lambda = \Sat_x(V_\lambda) \star \WW_\vac,\]
viewed a sheaf on $\Bun_N^{\lambda \cdot x}$. Since $\id \to V_\lambda^\vee \otimes V_\lambda$ is a direct summand, it suffices by $\Sph_{G,x}$-linearity to show that the cofiber of
    \[\Sat_x(V_\lambda) \star \eta_G \simeq \eta_{G,\Ran}(\WW_\lambda)\]
is cuspidal. In other words, we need to show for every proper parabolic $P$ that the morphism
\begin{equation}
\label{eqn:CT_eta}
    \CT_{P,!}^- \cdot \Poinc_{G,!} \cdot \WW_\lambda \to \CT_{P,!}^- \cdot \PsId_{G,!} \cdot \coeff_*^\vee \cdot \WW_{\lambda, \co}
\end{equation}
is an isomorphism.

\subsubsection{Third reduction} Let us rewrite the source of (\ref{eqn:CT_eta}) by use of the Jacquet functor
    \[
    \Jac_x : \Whit(G)_x \overset{\Jac_\Conf}{\longrightarrow} \Whit(M)_{\Conf_x} \overset{\sprd_*}{\longrightarrow} \Whit(M)_{\Ran_x}
    \]
and the coJacquet functor
    \[
    \co\Jac_x : \Whit(M)_{\Ran_x} \overset{\sprd^!}{\longrightarrow} \Whit(M)_{\Conf_x} \overset{\Jac_{\Conf,\co}}{\longrightarrow} \Whit(G)_x.
    \]
We obtain
\begin{equation}
\label{eqn:3_morphism_1}
\begin{split}
    & \CT_{P,!}^- \cdot \Poinc_{G,x,!} \cdot \WW_\lambda  \\
    & \qquad\simto \Poinc_{M,\Ran_x,!} \cdot \Jac_x \cdot \WW_\lambda \\
    & \qquad\qquad \simto \PsId_{M,!} \cdot \coeff_{M,\Ran_x,*}^\vee \cdot \left(\Jac_x \cdot \WW_\lambda \right)_\co \, [2\delta_M]\\
    & \qquad\qquad\qquad \simto \PsId_{M,!} \cdot \coeff_{M,\Ran_x,*}^\vee \cdot \co\Jac_x^\vee \cdot \WW_{\lambda, \co} \, [2\delta_G]
\end{split}
\end{equation}
The second arrow is induced the map $\eta_{M,\Ran_x}$, which is an isomorphism by induction. On the other hand, by rewriting the target of \ref{eqn:CT_eta}, we obtain
\begin{equation}
\label{eqn:3_morphism_2}
\begin{split}
    & \CT_{P,!}^- \cdot \Poinc_{G,x,!} \cdot \WW_\lambda \\
    & \qquad \to \CT_{P,!}^- \cdot \PsId_{G,!} \cdot \coeff_{G,x,*}^\vee \cdot \WW_{\lambda, \co} \, [2\delta_G] \\
    & \qquad\qquad \simto \PsId_{M,!} \cdot (\Eis_{P,!}^-)^\vee \cdot \coeff_{G,x,*}^\vee \cdot \WW_{\lambda, \co} \, [2\delta_G]\\
    & \qquad\qquad\qquad \simto \PsId_{M,!} \cdot \coeff_{M,\Ran_x, *}^\vee \cdot \Jac_{x,\co}^\vee \cdot \WW_{\lambda, \co} \, [2\delta_G].
\end{split}
\end{equation}
Therefore, it remains to perform the following:

\begin{construction}
\label{construction:3_morphism}
There is an identification of the morphism $(\ref{eqn:3_morphism_1})$ with $(\ref{eqn:3_morphism_2})$.
\end{construction}

\subsection{Construction of the identification}

\subsubsection{} We will compare (\ref{eqn:3_morphism_1}) with (\ref{eqn:3_morphism_2}) by describing both in explicit geometric terms. Consider the open substack
    \[
    \left( \Bun_N^{\lambda \cdot x} \underset{\Bun_G}{\times} \Bunt_{P^-} \underset{\Bun_M}{\times} \Bunc_M \right)^\trans
    \mon \Bun_N^{\lambda \cdot x} \underset{\Bun_G}{\times} \Bunt_{P^-} \underset{\Bun_M}{\times} \Bunc_M
    \]
cut out by the conditions:
\begin{itemize}
    \item the $N$ and $P^-$ structures are generically transverse, and
    \item the induced generic $N(M)$ structure is transverse to the $\Bunc_M$-structure.
\end{itemize}
This stack is equipped with with a natural projection $\pi$ onto $\Bun_M$.

\subsubsection{Description of (\ref{eqn:3_morphism_1})} Using the fact that $\WW_\lambda$ is a clean extension along the inclusion $\Bun_N^{\lambda \cdot x} \mon \Bun_N^{\pol\,x}$, it follows from a routine diagram chase that (\ref{eqn:3_morphism_1}) identifies with
    \[
    \pi_! \left( \WW_\lambda \stoximes j_{P,!}^- \stoximes r_{M,!} \right)^\trans 
    \to \pi_* \left( \WW_\lambda \shoximes j_{P,!}^- \shoximes r_{M,!} \right)^\trans [2 \delta_G + 2 \delta_M].
    \]
We shall identify (\ref{eqn:3_morphism_2b}) with this morphism.

\subsubsection{Functors defined by kernels} To handle (\ref{eqn:3_morphism_2}), we re-interpret the morphism
    \[
    \eta_G : \Poinc_{G,!} \cdot \WW_\lambda 
    \to \PsId_{G,!} \cdot \coeff_{x, *}^\vee \cdot \WW_{\lambda, \co} \, [2 \delta_G]
    \]
in terms of functors defined by kernels. Namely, we view
    \[
    \coeff_{\lambda, *} = \Gamma\left( \Bun_N^{\lambda \cdot x}, \WW_\lambda^{-1} \shimes \pi^!(-) \right)
    \qquad
    \coeff_{\lambda, !} = \Gamma_c\left(\Bun_N^{\lambda \cdot x}, \WW_\lambda^{-1} \stimes \pi^*(-) \right)
    \]
as functors defined by kernels
    \[
    \Sh(\Bun_G) \to \Sh(\pt).
    \]
By the assumption that $\Bun_N^{\lambda \cdot x}$ is a scheme, there is a canonical morphism
    \[\coeff_{\lambda, !} \to \coeff_{\lambda, *} \, [2 \delta_G].\]
Then $\eta_G$ identifies with the morphism
    \[
    (\id \boxtimes \coeff_{\lambda,!}) \cdot \Delta_{G,!} 
    \to (\id \boxtimes \coeff_{\lambda, *}) \cdot \Delta_{G,!} \, [2 \delta_G].
    \]

\subsubsection{Description of (\ref{eqn:3_morphism_2})} Now, factor (\ref{eqn:3_morphism_2}) as the composite of
\begin{equation}
\label{eqn:3_morphism_2a}
\begin{split}
    & \CT_{P,!}^- \cdot ( \Poinc_{G,!} \cdot \WW_\lambda ) \\
    & \qquad \to \CT_{P,!}^- \cdot ( \PsId_{G,!} \cdot \coeff_{x,*}^\vee \cdot \WW_{\lambda, \co} ) \, [2\delta_G] \\
    & \qquad\qquad \simto (\CT_{P,!}^- \cdot \PsId_{G,!}) \cdot \coeff_{x,*}^\vee \cdot \WW_{\lambda, \co} \, [2\delta_G] \\
    & \qquad\qquad\qquad \simto \PsId_{M,!} \cdot (\Eis_{P,!}^-)^\vee \cdot \coeff_{x,*}^\vee \cdot \WW_{\lambda, \co} \, [2\delta_G]
\end{split}
\end{equation}
and
\begin{equation}
\label{eqn:3_morphism_2b}
\begin{split}
    & \PsId_{M,!} \cdot (\Eis_{P,!}^-)^\vee \cdot \coeff_{x,*}^\vee \cdot \WW_{\lambda, \co} \, [2\delta_G] \\
    & \qquad \simto \PsId_{M,!} \cdot \coeff_{\Ran_x, *}^\vee \cdot \co\Jac_x^\vee \cdot \WW_{\lambda, \co} \, [2\delta_G]
\end{split}
\end{equation}
Then (\ref{eqn:3_morphism_2a}) identifies with the upper circuit of the diagram
\[
\begin{tikzcd}[
  column sep={18em,between origins},
  row sep={4em,between origins}
]
    (\CT_{P,!} \boxtimes \id) \cdot (\id \boxtimes \coeff_{\lambda,!}) \cdot \Delta_{G,!} 
    \arrow[r] \arrow[d] & (\CT_{P,!} \boxtimes \id) \cdot (\id \boxtimes \coeff_{\lambda, *}) \cdot \Delta_{G,!} \, [2 \delta_G] \arrow[d] \\
    (\id \boxtimes \coeff_{\lambda,!}) \cdot (\CT_{P,!} \boxtimes \id) \cdot \Delta_{G,!} 
    \arrow[r] \arrow[d] & (\id \boxtimes \coeff_{\lambda, *}) \cdot (\CT_{P,!} \boxtimes \id) \cdot \Delta_{G,!} \, [2 \delta_G] \arrow[d] \\
    (\id \boxtimes \coeff_{\lambda,!}) \cdot (\id \boxtimes \Eis_{P,!}^-) \cdot \Delta_{M,!}
    \arrow[r] & (\id \boxtimes \coeff_{\lambda, *}) \cdot (\id \boxtimes \Eis_{P,!}^-) \cdot \Delta_{M,!} \, [2 \delta_G]
\end{tikzcd}
\]
The lower face of this diagram is tautologically commutative. The commutation of the upper face is formal but tedious. Now reading the lower circuit of this diagram identifies (\ref{eqn:3_morphism_2a}) with
    \[
    \pi_! \left\{ \WW_\lambda \stoximes \left( j_{P,!}^- \stoximes r_{M,!} \right) \right\}^\trans 
    \to \pi_* \left\{ \WW_\lambda \stoximes \left( j_{P,!}^- \shoximes r_{M,!} \right) \right\}^\trans [2 \delta_G].
    \]
A further diagram chase identifies (\ref{eqn:3_morphism_2b}) with
    \[
    \pi_! \left\{ \WW_\lambda \stoximes \left( j_{P,!}^- \shoximes r_{M,!} \right) \right\}^\trans [2 \delta_G]
    \to \pi_* \left\{ \WW_\lambda \shoximes \left( j_{P,!}^- \shoximes r_{M,!} \right) \right\}^\trans [2 \delta_G + 2 \delta_M].
    \]
Composing these identifications completes Construction \ref{construction:3_morphism} and hence the proof of the main theorem.